\documentclass[../../main.tex]{subfiles}% 注意这里的文件路径不能用 ./main.tex ,否则用latexmk编译子文件会报错
\graphicspath{{\subfix{./image/}}{\subfix{../../image/}}} % 指定图片引用路径,后续可以直接使用图片文件名
% 如果图片存放在上级目录,则使用 ../../image/ 作为备用路径

% 例如:
% \begin{figure}[H]
% \centering
% \includegraphics[width=0.3\textwidth]{图.png}
% \caption{}
% \label{figure:图}
% \end{figure}
% 注意:上述\label{}一定要放在\caption{}之后,否则引用图片序号会只会显示??.

\begin{document}

\section{拓扑空间}

\begin{definition}
设 $X$ 是一个拓扑空间.
\begin{enumerate}[(1)]
\item 若 $X$ 不能表示为两个非空不相交的开集之并，则称 $X$ 是\textbf{连通的}.
\item 若 $A\subset X$ 作为子空间是连通的，则称 $A$ 为 $X$ 中的\textbf{连通子集}.
\end{enumerate}
\end{definition}

\begin{definition}
设 $X$ 是一个拓扑空间. 若对任意 $x,y\in X$，存在连续映射 $\gamma:[0,1]\to X$ 使得 $\gamma(0)=x$，$\gamma(1)=y$，则称 $X$ 是\textbf{道路连通的}. 子集 $A\subset X$ 若作为子空间道路连通，则称为 $X$ 中的\textbf{道路连通子集}.
\end{definition}

\begin{proposition}\label{proposition:道路连通空间一定是连通空间}
道路连通空间一定是连通空间.
\end{proposition}
\begin{proof}
设 $X$ 道路连通. 反设 $X$ 不连通，则可写成 $X=U\cup V$，其中 $U,V$ 为非空不相交开集. 取 $x\in U$，$y\in V$，由道路连通性，存在连续 $\gamma:[0,1]\to X$ 连接 $x$ 与 $y$. 则 $\gamma([0,1])$ 连通，但 $\gamma([0,1])\cap U$ 与 $\gamma([0,1])\cap V$ 是非空不相交的相对开集，其并集为 $\gamma([0,1])$，矛盾. 故 $X$ 连通.
\end{proof}

\begin{proposition}\label{proposition:道路连通空间一定是连通空间-1}
设 $\Omega\subset \mathbb{R}^n$ 是非空开集. 若 $\Omega$ 连通，则 $\Omega$ 道路连通.
\end{proposition}
\begin{remark}
由\refpro{proposition:道路连通空间一定是连通空间}和\refpro{proposition:道路连通空间一定是连通空间-1}知，$\mathbb{R}^n$ 中的非空开集连通当且仅当道路连通.
\end{remark}
\begin{proof}
固定 $x_0\in \Omega$. 定义
\begin{align*}
A \triangleq \{\, x\in \Omega \mid \text{存在道路在 }\Omega\text{ 内连接 }x_0\text{ 与 }x\,\}.
\end{align*}
显然 $A$ 非空（$x_0\in A$）. 下面证明 $A$ 既是 $\Omega$ 的开集又是闭集.

先证 $A$ 是开集：任取 $x\in A$，因 $\Omega$ 开，存在 $r>0$ 使得 $B(x,r)\subset \Omega$. 对任意 $y\in B(x,r)$，由于球 $B(x,r)$ 是凸集，线段连接 $x$ 与 $y$ 完全位于 $B(x,r)\subset \Omega$，故可从 $x_0$ 经道路到 $x$，再沿线段到 $y$，得 $y\in A$. 所以 $B(x,r)\subset A$，$A$ 为开集.

再证 $A$ 是闭集（在 $\Omega$ 中）：设 $z\in \Omega\setminus A$. 因 $\Omega$ 开，取 $r>0$ 使 $B(z,r)\subset \Omega$. 若存在 $y\in B(z,r)\cap A$，则从 $x_0$ 到 $y$ 有道路，再从 $y$ 到 $z$ 沿线段（在球内），可得 $z\in A$，矛盾. 故 $B(z,r)\cap A=\varnothing$，即 $\Omega\setminus A$ 为开集，所以 $A$ 在 $\Omega$ 中为闭集.

因 $\Omega$ 连通，且 $A$ 非空、既开又闭，故 $A=\Omega$. 于是 $\Omega$ 道路连通.
\end{proof}

\begin{definition}
设 $X$ 是一个拓扑空间.
\begin{enumerate}[(1)]
\item 若 $C\subset X$ 是连通的，且不存在 $X$ 中的连通子集 $D$ 使得 $C\subsetneq D$，则称 $C$ 为 $X$ 的一个\textbf{极大连通子集}.
\item 等价地，在 $X$ 上定义关系：$x\sim y$ 当且仅当存在一个包含 $x$ 与 $y$ 的连通子集. 该关系的每个等价类称为 $X$ 的一个\textbf{连通分支}.
\end{enumerate}
\end{definition}

\begin{proposition}\label{proposition:连通分支的基本性质}
设 $X$ 为拓扑空间,则
\begin{enumerate}[(1)]
\item\label{proposition:连通分支的基本性质-1} 连通分支一定是连通的.
\item\label{proposition:连通分支的基本性质-2} $E\subset X$是连通分支等价于$E$是$X$ 的极大连通子集；
\item\label{proposition:连通分支的基本性质-3} 不同连通分支不相交；
\item\label{proposition:连通分支的基本性质-4} 所有连通分支的并等于 $X$.
\end{enumerate}
\end{proposition}
\begin{proof}
\begin{enumerate}[(1)]
\item 由连通分支的定义知,每个等价类都是连通的,故任意两点可由连通子集相连，进而整个集合连通. 因此结论成立.

\item 先证每个等价类都是极大连通子集. 设 $E$ 是一个等价类. 任取 $x\in E$，因为 $x\sim x$，且单点集连通，所以 $E$ 非空. 对任意 $y,z\in E$，有 $x\sim y$ 且 $x\sim z$，因此存在连通集 $A$ 含 $x,y$，连通集 $B$ 含 $x,z$. 由于 $A\cup B$ 连通且含 $y,z$，故 $y\sim z$，从而 $E$ 是连通的（任意两点由连通子集相连）.

若 $E$ 不是极大的，则存在连通子集 $F$ 满足 $E\subsetneq F$. 取 $f\in F\setminus E$，任取 $e\in E$，则 $F$ 是包含 $e$ 与 $f$ 的连通集，故 $e\sim f$，与 $f\notin E$ 矛盾. 所以 $E$ 是极大连通子集，即连通分支.

反之，设 $C$ 是极大连通子集. 任取 $x\in C$，则 $C$ 是包含 $x$ 与 $y$ 的连通集，故对任意 $y\in C$ 有 $x\sim y$，即 $C$ 含于 $x$ 的等价类 $E_x$. 而 $E_x$ 连通（如上所证），且 $C\subset E_x$，由 $C$ 的极大性得 $C=E_x$. 因此每个极大连通子集恰为一个等价类.

\item 若两个连通分支 $C_1,C_2$ 相交，则 $C_1\cup C_2$ 是连通子集（两个连通集相交，并集连通），且 $C_1\subsetneq C_1\cup C_2$，这与 $C_1$ 的极大性矛盾.
\item 对任意 $x\in X$，由等价关系知 $x$ 属于某个等价类，即属于某个连通分支，因此并集为 $X$.
\end{enumerate}
\end{proof}

\begin{proposition}\label{proposition:Rn中开集的连通分支必是开集}
设 $\Omega\subset \mathbb{R}^n$ 是开集，$C$ 是 $\Omega$ 的一个连通分支，则
\begin{enumerate}[(1)]
\item\label{proposition:Rn中开集的连通分支必是开集-1} $C$是道路连通的.

\item\label{proposition:Rn中开集的连通分支必是开集-2} $C$是 $\mathbb{R}^n$ 中的开集.

\item\label{proposition:Rn中开集的连通分支必是开集-1-3} 在子空间拓扑下，$C$ 既是 $\Omega$ 的开子集也是 $\Omega$ 的闭子集.
\end{enumerate}
\end{proposition}
\begin{proof}
\begin{enumerate}[(1)]
\item 由\refpro{proposition:道路连通空间一定是连通空间-1}和\rrefpro{proposition:连通分支的基本性质}{proposition:连通分支的基本性质-1}立得.

\item 设 $C$ 是 $\Omega$ 的一个连通分支，任取 $x\in C$. 由于 $\Omega$ 是开集，存在 $r>0$ 使得 $B(x,r)\subset \Omega$. 因为 $B(x,r)$ 是连通的且与 $C$ 相交于 $x$，所以 $C\cup B(x,r)$ 是连通的. 由 $C$ 的极大性得 $C\cup B(x,r)\subset C$，从而 $B(x,r)\subset C$. 故 $C$ 是开集.

\item $C$ 是开集已由(1)得到. 又因 $\Omega\setminus C$ 是其余连通分支的并集，每个分支都是开集，故 $\Omega\setminus C$ 是开集，从而 $C$ 在 $\Omega$ 中是闭集.
\end{enumerate}
\end{proof}

\begin{proposition}\label{proposition:Rn中连通分支的边界包含于原集合的边界}
设 $\Omega\subset \mathbb{R}^n$ 是开集，$C$ 是 $\Omega$ 的一个连通分支，则
\begin{align*}
\partial C \subset \partial \Omega .
\end{align*}
即连通分支的边界包含于原集合的边界.
\end{proposition}
\begin{remark}
对 $\mathbb{R}^n$ 中的开集，连通性与道路连通性等价，因此同一连通分支内的任意两点都可由连续曲线连接.
\end{remark}
\begin{proof}
因为 $C\subset \Omega$，取闭包得 $\overline{C}\subset \overline{\Omega}$，故 $\partial C\subset \overline{C}\subset \overline{\Omega}$. 下面证明 $\partial C$ 不含 $\Omega$ 的内点.
假设存在 $x\in \partial C\cap \Omega$. 由\refpro{proposition:Rn中开集的连通分支必是开集}知$\Omega$ 是开集且 $C$ 是开集，因为开集不含边界点,所以$x\notin C$. 因此 $x$ 属于 $\Omega\setminus C$. 又因为连通分支构成 $\Omega$ 的划分，$x$ 必属于另一个连通分支 $C'$. 由于 $C'$ 也是开集，存在 $r>0$ 使得 $B(x,r)\subset C'$. 于是 $B(x,r)\cap C=\varnothing$，这与 $x\in \partial C$ 矛盾!故 $\partial C\cap \Omega=\varnothing$，结合 $\partial C\subset \overline{\Omega}$ 得 $\partial C\subset \overline{\Omega}\setminus \Omega=\partial \Omega$.
\end{proof}

\begin{theorem}\label{theorem:紧拓扑空间上的连续实值函数必有最值}
设$X$是紧拓扑空间,$f:X\to \mathbb{R}$连续,则$f$在$X$上必能取到最大值和最小值.
\end{theorem}
\begin{proof}


\end{proof}





















\end{document}